\usepackage{listings}
\usepackage{xcolor}
\usepackage{float}
\usepackage{svg}
\usepackage{siunitx}
\usepackage{caption}
% Define colors for syntax highlighting
\definecolor{codegreen}{rgb}{0,0.6,0}
\definecolor{codegray}{rgb}{0.5,0.5,0.5}
\definecolor{codepurple}{rgb}{0.58,0,0.82}
\definecolor{backcolour}{rgb}{0.95,0.95,0.92}

% Configure listings style
\lstdefinestyle{mystyle}{
    backgroundcolor=\color{backcolour},
    commentstyle=\color{codegreen},
    keywordstyle=\color{magenta},
    numberstyle=\tiny\color{codegray},
    stringstyle=\color{codepurple},
    basicstyle=\ttfamily\footnotesize,
    breakatwhitespace=false,
    breaklines=true,
    captionpos=b,
    keepspaces=true,
    numbers=left,
    numbersep=5pt,
    showspaces=false,
    showstringspaces=false,
    showtabs=false,
    tabsize=2
}
\lstset{style=mystyle}

%% include \begin{comment} and \end{comment}
\usepackage{verbatim}

%% Set paper and border size
\usepackage[paper=a4paper,left=25mm,right=25mm,top=25mm,bottom=25mm]{geometry}

%% UTF8 as input characters; UCS incompatible to biblatex
\usepackage[utf8]{inputenc}

%% The default setting of the language is American. Please changesettings for
%% additional or alternative languages used in main.tex.
%% Please note that the default language of the document is the *last*language
%% which is added to the package options.
%% To set only parts of your document in a different language as the rest,
%% use for example\newline\verb+\foreignlanguage{ngerman}{Beispieltext indeutscher Sprache}+\newline
\usepackage[\mylanguage]{babel}

%% This package defines basic colors. If you want to get rid of colored links
%% and headings please change corresponding value in main.tex to {0,0,0}.
%% Used for links and so forth in screen-version
\usepackage[dvipsnames]{xcolor}
\definecolor{DispositionColor}{RGB}{\mydispositioncolor}

%% The widely used package to use graphical images within a LaTeXdocument.
%% \includegraphics[width=42mm]{figures/image}
\usepackage{graphicx}

%% For example on title pages you might want to have a logo on the upperright
%% corner of the first page (only). The package \texttt{eso-pic} is ableto
%% place things on absolute and relative positions on the whole page.
\usepackage{eso-pic}

\usepackage{tikz}
\usetikzlibrary{shapes.geometric, arrows.meta, positioning}
\usepackage{booktabs}

%% List of abbreviations
\usepackage{acronym}

   %% Use biblatex for bibliography
   %% With "Sorting=None" the numbering is done according to the order of appearance. 
   %\usepackage[style=numeric,sorting=none]{biblatex}
   
   %% Used for quotes
   \usepackage{csquotes}
   
   %% Add fontec to hyphenate umlaut and fix-cm to fix changes in section heading
   \usepackage[T1]{fontenc}
   \usepackage{fix-cm}

   %% This document template is able to generate an output that uses colorized
   %% headings, captions, page numbers, and links. The color named
   %% `DispositionColor' used in this document is defined near the definition
   %% of package xcolor
   %% The changes required for headings, page numbers and captions are defined
   %% here.
   %% Settings for colored links are handled by the definitions of the
   %% hyperref package
   %% \setheadsepline{.4pt}[\color{DispositionColor}]
   %\renewcommand{\headfont}{\normalfont\sffamily\color{DispositionColor}}
   %\renewcommand{\pnumfont}{\normalfont\sffamily\color{DispositionColor}}
   %\addtokomafont{disposition}{\color{DispositionColor}}
   %\addtokomafont{caption}{\color{DispositionColor}\footnotesize}
   %\addtokomafont{captionlabel}{\color{DispositionColor}}


    \renewcommand{\familydefault}{\sfdefault}

   %% Add package to add long text tables
   \usepackage{longtable}
   
   \usepackage{etoolbox}
   
   %% for the code-block environment with listing and reference possibility:
   %% \begin{code}{<language>}{<caption>}{<label>} 
   %% ... 
   %% \end{code}
   %% is also used with listings so there is a list of code blocks
   \usepackage{listings}
   \usepackage[figure,table,lstlisting]{totalcount}
   %\renewcommand{\lstlistingname}{Code}
   \usepackage{caption}

   %% Use caption package to customize image, code and table captions
   \usepackage[
       font =      {footnotesize, sf}, %% Sans Serif font family
       labelfont = {footnotesize, bf} %% Bold font series
   ]{caption}
   
   %% include amsmath to include text like ä,ö,ü,ß in equation
   \usepackage{amsmath}
   
   %% set bibliographystyle to DHBW appropriate
   \bibliographystyle{abbrvdin}
   
   %% to quote directly in text
   \newcommand*{\zitext}[2]{%
      \glqq#1\grqq\ %
      #2%
   }

   %% include package sinunitx to format units
   \usepackage[ 
   locale=DE, 
   separate-uncertainty, 
   mode = text,
   detect-all,
   multi-part-units      = single, 
   range-units           = single, 
   list-units           = single, 
   ]{siunitx}
   
   \DeclareSIUnit\fps{fps}    %% declare own units e.g. fps
   
   % to include pdfs as img or pdfs
   \usepackage{graphics}
   \usepackage{pdfpages}
   
   % rotating of e.g. tables
   \usepackage{rotating}
   
   % combining of rows in tables
   \usepackage{multirow}
   
   \usepackage{subcaption}
   
   %\usepackage{graphicx}
   
   \usepackage{etoolbox}
   
   % set table font to sf
   \AtBeginEnvironment{tabular}{\sffamily}
   
   \usepackage{listings}
   
   %\lstdefinestyle{mystyle}{
   %    basicstyle=\sffamily\small,
   %    
   %}
   %
   %\lstset{style=mystyle}

%\usepackage[hidelinks]{hyperref}